\section{数据结构}

\subsection{二叉树基础}
\subsubsection{先序遍历}
\begin{lstlisting}[language=c++,caption={先序遍历}]
#include <bits/stdc++.h>
using namespace std;

typedef struct Tree {
	struct Tree *left;
	struct Tree *right;
	int data;
} Tree, *BitTree;

void PreOrder(BitTree BT) {
	if (BT) {
		cout << BT->data << " ";
		PreOrder(BT->left);
		PreOrder(BT->right);
	}
}
\end{lstlisting}

\subsubsection{中序遍历}

\begin{lstlisting}[language=c++,caption={中序遍历}]
#include <bits/stdc++.h>
using namespace std;

typedef struct Tree {
	struct Tree *left;
	struct Tree *right;
	int data;
} Tree, *BitTree;

void PostOrder(BitTree BT) {
	if (BT) {
		PostOrder(BT->left);
		cout << BT->data << " ";
		PostOrder(BT->right);
	}
}
\end{lstlisting}


\subsubsection{后序遍历}
\begin{lstlisting}[language=c++,caption={后序遍历}]
#include <bits/stdc++.h>
using namespace std;

typedef struct Tree {
	struct Tree *left;
	struct Tree *right;
	int data;
} Tree, *BitTree;

void PostOrder(BitTree BT) {
	if (BT) {
		PostOrder(BT->left);
		PostOrder(BT->right);
		cout << BT->data << " ";
	}
}
\end{lstlisting}

\subsubsection{并查集}


\subsection{线段树模板}

\begin{lstlisting}[language=c++,caption={线段树模板}]
//线段树模板
#include<bits/stdc++.h>
using namespace std;
#define lc p<<1
#define rc p<<1|1
#define N 500005
#define int long long
int m, n, w[N];

struct node {
	int l, r, sum, add;
} tr[N * 4];

void build(int p, int l, int r) {
	tr[p] = {l, r, w[l], 0};
	if (l == r) {
		return;
	}
	int m = (l + r) >> 1;
	build(lc, l, m);
	build(rc, m + 1, r);
	tr[p].sum = tr[lc].sum + tr[rc].sum;
}

void pushup(int p) {
	tr[p].sum = tr[lc].sum + tr[rc].sum;
}

void pushdown(int p) {
	if (tr[p].add) {
		tr[lc].sum += tr[p].add * (tr[lc].r - tr[lc].l + 1);
		tr[rc].sum += tr[p].add * (tr[rc].r - tr[rc].l + 1);
		tr[lc].add += tr[p].add;
		tr[rc].add += tr[p].add;
		tr[p].add = 0;
	}
}

void update(int p, int x, int y, int k) {
	if (x <= tr[p].l && tr[p].r <= y) {
		tr[p].sum += (tr[p].r - tr[p].l + 1) * k;
		tr[p].add += k;
		return;
	}
	int m = (tr[p].l + tr[p].r) >> 1;
	pushdown(p);
	if (x <= m) {
		update(lc, x, y, k);
	}
	if (y > m) {
		update(rc, x, y, k);
	}
	pushup(p);
}

int query(int p, int x, int y) {
	if (x <= tr[p].l && tr[p].r <= y) {
		return tr[p].sum;
	}
	int m = (tr[p].l + tr[p].r) >> 1;
	int sum = 0;
	pushdown(p);
	if (x <= m) sum += query(lc, x, y);
	if (y > m) sum += query(rc, x, y);
	return sum;
}

void solve() {
	cin >> n >> m;
	for (int i = 1; i <= n; i++) {
		cin >> w[i];
	}
	build(1, 1, n);
	while (m--) {
		int op, x, y, k;
		cin >> op;
		switch (op) {
			case 1:
				cin >> x >> y >> k;
				update(1, x, y, k);
				break;
			case 2:
				cin >> x >> y;
				cout << query(1, x, y) << endl;
			default:
				break;
		}
	}
}

signed main() {
	ios_base::sync_with_stdio(false);
	cin.tie(NULL);
	int t = 1;
	while (t--) {
		solve();
	}
	return 0;
}

\end{lstlisting}


\subsection{树状数组模板}

\begin{lstlisting}[language=c++,caption={树状数组模板}]
//树状数组模板
#include<bits/stdc++.h>
using namespace std;
#define lc p<<1
#define rc p<<1|1
#define N 500005
#define int long long
int m, n, w[N];

struct node {
	int l, r, sum, add;
} tr[N * 4];

void build(int p, int l, int r) {
	tr[p] = {l, r, w[l], 0};
	if (l == r) {
		return;
	}
	int m = (l + r) >> 1;
	build(lc, l, m);
	build(rc, m + 1, r);
	tr[p].sum = tr[lc].sum + tr[rc].sum;
}

void update(int p, int x, int k) {
	if (tr[p].l == x && tr[p].r == x) {
		tr[p].sum += k;
		return;
	}
	int m = (tr[p].l + tr[p].r) >> 1;
	if (x <= m) update(lc, x, k);
	if (x > m) update(rc, x, k);
	tr[p].sum = tr[lc].sum + tr[rc].sum;
}

void pushup(int p) {
	tr[p].sum = tr[lc].sum + tr[rc].sum;
}

void pushdown(int p) {
	if (tr[p].add) {
		tr[lc].sum += tr[p].add * (tr[lc].r - tr[lc].l + 1);
		tr[rc].sum += tr[p].add * (tr[rc].r - tr[rc].l + 1);
		tr[lc].add += tr[p].add;
		tr[rc].add += tr[p].add;
		tr[p].add = 0;
	}
}

int query(int p, int x, int y) {
	if (x <= tr[p].l && tr[p].r <= y) {
		return tr[p].sum;
	}
	int m = (tr[p].l + tr[p].r) >> 1;
	int sum = 0;
	pushdown(p);
	if (x <= m) sum += query(lc, x, y);
	if (y > m) sum += query(rc, x, y);
	return sum;
}

void solve() {
	cin >> n >> m;
	for (int i = 1; i <= n; i++) {
		cin >> w[i];
	}
	build(1, 1, n);
	while (m--) {
		int op, x, y, k;
		cin >> op;
		switch (op) {
			case 1:
				cin >> x >> k;
				update(1, x, k);
				break;
			case 2:
				cin >> x >> y;
				cout << query(1, x, y) << endl;
			default:
				break;
		}
	}
}

signed main() {
	ios_base::sync_with_stdio(false);
	cin.tie(NULL);
	int t = 1;
	while (t--) {
		solve();
	}
	return 0;
}	
\end{lstlisting}

\subsection{平衡树}




